\begin{ozet}
Bu tez, bir Android telefon üzerinde çalışan uçtan uca bir
Amerikan İşaret Dili (ASL) tanıma sistemi sunar ve iki görevi ele
alır: ASL Alphabet veri kümesinin 29 sınıfı üzerinde harf düzeyinde
parmak alfabesi tanıma ve WLASL-100 veri kümesinin 100 kelimelik
sözdağarcığı üzerinde kelime düzeyinde tanıma. Harf görevi için iki
sınıflandırıcı karşılaştırılır: ham görüntüleri işleyen bir
EfficientNet-B0 ve MediaPipe Hands tarafından çıkarılan 21 el işaret
noktası üzerinde çalışan kompakt bir çok katmanlı algılayıcı (MLP).
Tek imzalayıcılı ASL Alphabet test bölümünde sırasıyla \%100 ve
\%99,38 Top-1 değerlerine ulaşırlar; bu sayılar bir yetenek iddiası
değil, eğitim hattının sağlık kontrolüdür; landmark model 0,24
MB'tir. Yeniden eğitilmeden, farklı imzalayıcı ve arka planlardan
oluşan bağımsız \num{870} görüntülük bir dış kümede görüntü ağının
doğruluğu \%59,77'ye çökerken landmark MLP tespit edilen ellerde
\%90,72'yi korur; bu, görüntü modelinin çekim-ortamı ezberini,
landmark modelinin ortamdan bağımsızlığını ortaya koyar. Cihaz üzeri
değerlendirme de bütçe telefonda bu sıralamayı tersine çevirir ve harf
görevinde landmark türevini varsayılan yapar.
Kelime görevi için, Kaggle paketindeki \num{1013} video yt-dlp ile
\num{1465} klibe çıkarılmış (\%44 artış) ve iki elden çıkarılan 128
boyutlu işaret noktası dizisi üzerinde çift yönlü bir LSTM
eğitilmiştir. Beş seed'li logit ortalamalı ensemble, person-mixed
bölümlemede \%56,22 Top-1 ve \%81,09 Top-5'e ulaşır ve tezin
benimsediği \%40 eşiğini 16 puan aşar; person-mixed olduğundan bu,
signer-bağımsız değil, literatürle kıyaslanabilir bir değerdir. Dokuz
adımlı bir ablasyon, \%20 tabanından \%60,08 doğrulamaya çıkan
kazanımı ağırlığı benzer üç katkıya ayırır: veri mühendisliği, çıkarım
hattı ve eğitim rejimi. Dört sınıflandırıcı TensorFlow Lite'a
aktarılıp CameraX ve MediaPipe Hands tabanlı bir Android
uygulamasında dağıtılmış; FPS, gecikme ve OOD farkı bütçe
(TECNO SPARK 20C) ve amiral gemisi (Xiaomi 15T Pro) cihazlarda
ölçülerek cihaza bağlı bir model seçim rehberi sunulmuştur.
\end{ozet}
